\section{Correctness of the recursive solver}
\label{sec:solving}

The orientation alternative now supplies the missing ingredient for
\(\Solve\): each fixed admissible subsolution lies below its entire
output with probability at least \(15/16\), and each fixed admissible
supersolution lies above it with the same probability. We prove these
two assertions separately by induction on \(d(R)\). The transformed
operator and all global parameters remain fixed.

The main continuation preserves rank, so its comparison guarantee is
not yet available in this induction. Instead, we follow one comparison
through the entire chain of main continuations. At each node, the
local calls either establish the required inequality or pass one
admissible comparison to the next node. Bounding the probability of
the first local error along this finite chain closes the induction.
Throughout the section, \(\varepsilon=2^{-J}\) is analysis notation.

\subsection{Using a pointwise guarantee after earlier random choices}

A transported comparison can depend on earlier outputs. What matters
is that the comparison and its admissibility are settled before the
fresh call that will test it. We record the conditioning rule explicitly.

\begin{lemma}[Comparisons determined by the past]
\label{sol:conditioning}
Suppose a randomized procedure has failure probability at most \(\rho\)
for every fixed input and every fixed comparison satisfying a specified
qualification condition. Let \(\mathcal H\) be the history before a
call. Assume that its input, proposed comparison, and qualification are
determined by \(\mathcal H\), and that the call uses fresh randomness.
Then
\[
 \Pr\bigl(\text{qualification holds and the comparison fails}
              \mid\mathcal H\bigr)
 \le \rho\,\mathbf 1_{\{\text{qualification holds}\}}.
\]
The same bound may be used after restricting to any event determined
by this pre-call history.
\end{lemma}

\begin{proof}
Condition on a possible history. If the comparison does not qualify,
the event on the left is empty. Otherwise that history specifies one
deterministic eligible input and comparison. The fresh bits retain
their original distribution independently of the history, so the
assumed pointwise bound applies. Averaging these conditional bounds
proves the assertion, including its restriction to a pre-call event.
\end{proof}

We will use the lemma in a first-error argument. Suppose that an
analysis requests at most \(M\) comparisons, ceasing its requests
after the first failure. Each request and its qualification must be
determined before the corresponding call. If every qualified request
has conditional failure probability at most \(\rho\), the probability
that a first failure occurs is at most \(M\rho\). Indeed, reaching a
specified request with no earlier failure is a pre-call event. Apply
Lemma~\ref{sol:conditioning} on that event and sum over the \(M\)
possible positions, treating unreached positions as empty events.
This argument requires only the single comparison at each position.

Under the induction hypothesis for smaller ranks,
Lemma~\ref{lem:ask-amplification} gives this rule for \(\Ask\) with
\(\rho=\varepsilon\). Proposition~\ref{prop:orientation} and
Lemma~\ref{lem:orient-plus} give it for \(\OrientPlus\) whenever the
fixed comparison satisfies the required core condition. Those results
use only smaller-rank solver guarantees; they are therefore available
before we prove correctness at the current rank.

\subsection{Preserving a comparison while tightening the box}

Tightening has two roles. It transports the comparison through each
change of bounds, and it ensures that a comparison threatened by the
subsequent truncation has a heavy extreme core. That second property
is precisely the hypothesis needed by the orientation alternative.

\begin{lemma}[Tightening a fixed comparison]
\label{sol:tightening}
Assume the comparison guarantees at every smaller rank. Consider a
\(\Solve\) tightening loop with \(B>128s\), and fix an admissible
subsolution \(y^{\mathrm{in}}\) before the loop begins. Outside an
event of probability at most \((16ns+2)\varepsilon\), its final box
\([a,b]\) contains an admissible subsolution
\(y\ge y^{\mathrm{in}}\), obtained by successive lower-bound
projections, such that
\begin{equation}
 \max_i(y_i-a_i)>C
 \quad\Longrightarrow\quad
 \alpha\bigl(\{i:y_i\ge a_i+B-hB\}\bigr)>\eta k,
 \qquad C=B-\delta.
 \label{sol:tightened-high-core}
\end{equation}
Separately, for each admissible supersolution
\(z^{\mathrm{in}}\) fixed before the loop, outside an event with the
same probability bound, the final box contains an admissible
\(z\le z^{\mathrm{in}}\) satisfying
\begin{equation}
 \max_i(b_i-z_i)>C
 \quad\Longrightarrow\quad
 \beta\bigl(\{i:z_i\le b_i-B+hB\}\bigr)>\eta m.
 \label{sol:tightened-low-core}
\end{equation}
Both assertions also hold conditionally on any preceding history that
determines the loop input and the initial comparison.
\end{lemma}

\begin{proof}
We specify at most one comparison to request from each actual
\(\Ask\) call. The analysis stops at its first failed request; the
algorithm continues to execute its prescribed instructions on every
tape.

\paragraph{The subsolution.}
Let \(y\) denote the current protected projection. Before a
reduced-\(k\) call, the box, \(y\), and
\(d=\max_i(y_i-a_i)\) are already determined. If \(d>C\), form
\[
 u=\max(y-d+\delta,a).
\]
By Lemma~\ref{lem:tightening-projection}, this is a subsolution of
the current box. Its support is contained in
\(I_y=\{i:y_i\ge a_i+B-hB\}\): at an active coordinate,
\[
 y_i-a_i>d-\delta>B-2\delta\ge B-hB.
\]
At a coordinate attaining \(d\), the projected value is
\(u_i=a_i+\delta\). If \(\alpha(S_u)\le\eta k\), request that
the call return a vector at least \(u\). The vector and this
qualification are fixed before the call, and success forces a
lower-bound update. If \(d\le C\), or if \(u\) fails to qualify,
make no request at this call.

Whenever the algorithm raises the lower bound to an output \(r\),
replace the protected vector by \(\max(y,r)\). Both vectors lie
below the unchanged upper bound. Lemma~\ref{lem:boundary-projections}
therefore gives an admissible subsolution for the new box, with no
larger support and no smaller coordinates. This projection needs only
the unconditional fact that \(r\) lies in the old box.

At every reduced-\(m\) call that is reached, request comparison with
the current \(y\). Its full subsolution weights and budget are
unchanged, so it qualifies. Success gives \(r\ge y\); lowering the
upper bound to \(r\), if an update occurs, still contains \(y\)
and leaves its defining support and inequalities unchanged.

Consider now the final pass, in which neither bound changes. Its
reduced-\(k\) call used the final box and the final protected vector.
If \(d>C\) but \(\alpha(I_y)\le\eta k\), the candidate \(u\)
at that call qualifies because \(S_u\subseteq I_y\). In the absence
of a charged failure its output forces an update, contradicting that
this was the final pass. This proves
Equation~\eqref{sol:tightened-high-core}. Every candidate was fixed
before its own call; identifying the final pass afterwards makes no
additional probabilistic demand.

\paragraph{The supersolution.}
Let \(z\) now denote the current protected vector. At each
reduced-\(k\) call request an output at most \(z\). The call
retains the full supersolution parameters, and its successful output
keeps \(z\) inside the box if the lower bound is raised.
Before a reduced-\(m\) call, put \(e=\max_i(b_i-z_i)\). If
\(e>C\), define
\[
 v=\min(z+e-\delta,b).
\]
Lemma~\ref{lem:tightening-projection} makes \(v\) a supersolution.
Its support lies in \(I_z=\{i:z_i\le b_i-B+hB\}\), since an
active coordinate satisfies
\(b_i-z_i>e-\delta>B-2\delta\ge B-hB\).
A coordinate attaining \(e\) has \(v_i=b_i-\delta\).
Whenever \(v\) qualifies for the reduced budget, request an output
at most \(v\); success forces an upper-bound update. After any
upper-bound update to \(r\), project \(z\) to \(\min(z,r)\).
Lemma~\ref{lem:boundary-projections} preserves admissibility and
keeps the support from increasing. At the final pass, \(e>C\)
together with \(\beta(I_z)\le\eta m\) would again force an update.
This proves Equation~\eqref{sol:tightened-low-core}.

\paragraph{The error count.}
For either fixed side, every requested comparison qualifies and is
determined before its call until the first error. There is at most
one request per call. Proposition~\ref{prop:termination} bounds the
loop by \(8ns\) updates, hence \(8ns+1\) passes and \(16ns+2\)
calls. The first-error bound from Lemma~\ref{sol:conditioning} gives
\((16ns+2)\varepsilon\). Conditioning first on the history before
the loop gives the asserted conditional version.
\end{proof}

\subsection{A comparison after the choice of boundary}

After successful tightening, fix the resulting comparison in the
final box \([a,b]\). The orientation call and all reweighted calls
use this same box. We next determine exactly which local comparisons
suffice at this node. In each case, either a local answer supplies
the desired inequality, or one admissible comparison remains for the
main successor. We assume no correctness guarantee for that successor.

\paragraph{A subsolution and the returned minimum or maximum.}
Let \(y\) be the current subsolution. When
\(\max_i(y_i-a_i)>C\), Equation~\eqref{sol:tightened-high-core}
allows us to request the \(\OrientPlus\) guarantee for \(y\).
This vector is fixed before the orientation randomness; the
conditional failure probability is at most \(\varepsilon\).
When \(\max_i(y_i-a_i)\le C\), no orientation guarantee is needed.

\begin{enumerate}[label=\textup{(\roman*)},leftmargin=*]
\item
Suppose the choice is \(Y\), so the node returns a minimum.
The reweighted calls change only the supersolution weights and
budget. The vector \(y\) therefore qualifies in every one of them;
request a lower comparison with \(y\) from each call. Define
\[
 a'=\max(a,b-C),\qquad y'=\max(y,a').
\]
By Lemma~\ref{lem:boundary-projections}, \(y'\) is an admissible
subsolution in the main box \([a',b]\), with the unchanged
subsolution parameters. If the main output is at least \(y'\) and
all requested local outputs are at least \(y\), their minimum is
at least \(y\). Pass the single comparison \(y'\) to the main
successor.

\item
Suppose the choice is \(Z\), so the node returns a maximum, and put
\(b'=\min(b,a+C)\). If \(y\le b'\), the same vector remains an
admissible subsolution in \([a,b']\). An output at least \(y\)
from the main successor suffices for the returned maximum. Pass this
comparison to that successor; no reweighted output need be tested.

If instead \(y\not\le b'\), then \(y\le b\) implies
\(\max_i(y_i-a_i)>C\). A successful requested orientation
guarantee supplies a listed set \(H\) satisfying
\begin{equation}
 \alpha(S_y\cap H)>\eta k.
 \label{sol:witness-high}
\end{equation}
For the analysis, select the first such entry. The reweighting lemma
makes \(y\) admissible for that entry's call. Request its lower
comparison with \(y\). Success makes the returned maximum at least
\(y\), whatever the main output may be. The obligation ends at
this node.
\end{enumerate}

The selected witness in Equation~\eqref{sol:witness-high} is a
function of the completed list and the vector fixed before orientation.
Thus its index is determined before any reweighted call begins.
The algorithm executes every list entry without computing this
analysis-dependent index. Earlier entries may already have returned
when the selected entry is reached, but its own bits are still fresh.
Lemma~\ref{sol:conditioning} therefore applies to a witness chosen
from the random list.

\paragraph{A supersolution and the returned maximum or minimum.}
For the current supersolution \(z\), request the orientation
guarantee exactly when \(\max_i(b_i-z_i)>C\); the needed core
condition is Equation~\eqref{sol:tightened-low-core}.
The two output rules give the following obligations.

\begin{enumerate}[label=\textup{(\roman*)},leftmargin=*]
\item
For choice \(Z\), all reweighted calls retain the supersolution
parameters. Request an upper comparison with \(z\) from every
entry. Put \(b'=\min(b,a+C)\) and \(z'=\min(z,b')\).
Lemma~\ref{lem:boundary-projections} makes \(z'\) admissible in
\([a,b']\). If the main output is at most \(z'\), and all
reweighted outputs are at most \(z\), their maximum is at most
\(z\). Pass \(z'\) to the main successor.

\item
For choice \(Y\), put \(a'=\max(a,b-C)\). If \(z\ge a'\),
it remains admissible in \([a',b]\). A main output at most \(z\)
already controls the returned minimum, so pass that comparison to
the main successor.

If \(z\not\ge a'\), the inequality \(z\ge a\) implies
\(\max_i(b_i-z_i)>C\). Successful orientation supplies an entry
\(H\) with \(\beta(S_z\cap H)>\eta m\). Select the first such
entry before its call. Lemma~\ref{lem:reweighting} makes \(z\)
admissible there. Its successful upper comparison forces the
returned minimum to be at most \(z\), completing the obligation
without a condition on the main output.
\end{enumerate}

The minimum for a subsolution under choice \(Y\), and the maximum
for a supersolution under choice \(Z\), require all the reweighted
comparisons just specified. The witness cases require only one.
In either situation, Proposition~\ref{prop:termination} bounds the
number of entries by \(J(2g+2)\) on every tape. Each entry has a
fixed input and comparison before its own fresh call. Applying the
first-error bound to these entries, together with the one possible
orientation request, charges at most
\begin{equation}
 \bigl(1+J(2g+2)\bigr)\varepsilon
 \label{sol:post-tightening-error}
\end{equation}
at this node for either protected side. The internal calls of
\(\OrientPlus\) are included in its guarantee and contribute no
second error charge here.

\subsection{Closing the rank induction along a main chain}

\begin{theorem}[Pointwise comparisons for the solver]
\label{thm:solve-comparisons}
Let \([a,b]\) be an allowed integer box with positive integer
reference width \(B\le B_*\), and let the support weights and
budgets be positive rationals. For the output \(p\) of \(\Solve\),
each fixed admissible subsolution \(y\) and, separately, each fixed
admissible supersolution \(z\) satisfy
\[
 \Pr(p\ge y)\ge 1-\frac1{16},
 \qquad
 \Pr(p\le z)\ge 1-\frac1{16}.
\]
These are whole-vector inequalities, as in
Equation~\eqref{eq:comparison-guarantees}: a fixed comparison is
chosen before the invocation, and its inequality holds at every
coordinate on its success event.
\end{theorem}

\begin{proof}
Induct on \(d(R)\), allowing every permitted box and all positive
rational support parameters at each rank. Rank zero means \(R<1\),
so Lemma~\ref{lem:small-rank} proves both assertions without error.
At any rank, the small-width branch \(B\le128s\) returns the exact
box solution. Lemmas~\ref{lem:box-comparison} and
\ref{lem:box-solution} put it above every subsolution and below
every supersolution.

Now fix a positive rank and assume the assertions at all smaller
ranks. Budget reductions and the reweightings of
Lemma~\ref{lem:reweighting} strictly decrease rank. Thus the
amplification and orientation results, and all conditional estimates
proved above, are available. In particular, orientation correctness
has already been established by its own width induction from these
smaller-rank solver guarantees.

Fix one admissible subsolution before the present invocation starts.
Follow its obligation along the actual main chain according to the
preceding subsection. On arrival at a node, the entry box and
comparison are functions of the computations already completed.
Until the first charged error, tightening preserves admissibility
and increases the protected vector. The subsequent local work either
completes its obligation or passes one admissible comparison to the
main successor. In the latter case, we analyze that successor's
local work next; we do not invoke an unproved same-rank guarantee.

By Equation~\eqref{eq:main-chain-length}, every tape has at most
\(K=8s(g+1)+1\) nodes in this chain. Combining the tightening count
with Equation~\eqref{sol:post-tightening-error}, the first-error
charge at any one reached node is at most
\(\Lambda\varepsilon\), where
\begin{equation}
 \Lambda=16ns+3+J(2g+2)
 \le 100s(n+1)+10J(g+2).
 \label{sol:node-error}
\end{equation}
The \(16ns+2\) possible tightening requests give the first part of
the count, followed by one orientation request and at most
\(J(2g+2)\) reweighted requests. This conditional charge applies
on precisely the histories that reach the node with no earlier
error, for those histories retain the necessary comparison
invariants. The first-error union bound over the whole chain gives
total probability at most \(K\Lambda\varepsilon\).

To verify success when no charged error occurs, start at the node
that completes the obligation, or at the exact terminal solver if
it is passed all the way to the end. Its output has the required
inequality. Work backwards through the preceding nodes. Each has
the local inequalities required by its maximum or minimum rule,
and the main successor supplies the one passed inequality.
Consequently its output dominates its protected vector. The
tightening projections dominate the vector at that node's entry,
so the inequality propagates to the initial comparison. The
algorithm still executes a main call below a node that completes
the obligation; that computation is irrelevant to this implication.

For a fixed supersolution, follow instead the upper comparisons
specified above. Tightening decreases the protected vector, the
passed comparisons remain admissible, and the same backwards
argument establishes the original upper bound. Its request count
is again \(\Lambda\). This is a separate analysis for the chosen
supersolution; no guarantee simultaneous over all admissible
comparison vectors has been used.

It remains to check the numerical error bound. Set \(Q=s(n+1)\).
The global choices give \(Q\ge2\), \(g+2\le s\le Q\), and
\(J=2Q^6+1\ge32\). In particular,
\[
 K\le 9Q^2,
 \qquad
 \Lambda\le100Q+10J(g+2)\le20JQ,
\]
where the last step uses \(J\ge10\). Since \(Q^3\le J\),
\begin{equation}
 K\Lambda\varepsilon
 \le180JQ^3\,2^{-J}
 \le180J^2\,2^{-J}
 <\frac1{16}.
 \label{sol:chain-error-bound}
\end{equation}
Indeed, \(J^2 2^{-J}\) decreases for integers \(J\ge3\), and
\(180\cdot32^2 2^{-32}<1/16\). This proves both pointwise
assertions at the current rank and completes the induction.
\end{proof}

\subsection{Recovering the entire top-box potential}

The two pointwise inequalities can now be applied to one deterministic
vector: the exact solution of the initial box. Their intersection
recovers every coordinate, so the success probability needs no union
bound over vertices.

\begin{corollary}[Exact recovery of the top box solution]
\label{cor:top-success}
Let \(x\) be the box solution for \([0,B_*]\). Let \(p\) be the
output of \(\Solve\) on that box with unit support weights,
budgets \(k=m=n\), and reference width \(B_*\). Then
\[
 \Pr(p=x)\ge\frac78.
\]
Hence \(\{i:p_i>B_*/2\}\) is the complete winning set with
probability at least \(7/8\).
\end{corollary}

\begin{proof}
The vector \(x\) is both a subsolution and a supersolution, and
each support has cardinality at most \(n\). It qualifies on both
sides for the prescribed unit weights and budgets. Apply
Theorem~\ref{thm:solve-comparisons} to this same deterministic
vector in its two roles and to the same output \(p\). Then
\[
 \Pr(p=x)
 =\Pr(p\ge x\text{ and }p\le x)
 \ge1-\Pr(p\not\ge x)-\Pr(p\not\le x)
 \ge1-\frac1{16}-\frac1{16}=\frac78.
\]
The intersection is equality of entire vectors; independence of the
two events is unnecessary. On that event,
Proposition~\ref{prop:game-reduction} identifies the thresholded
output with the complete winning set.
\end{proof}
